\documentclass[../main.tex]{subfiles}
\begin{document}

\chapter{Estado del arte}
\label{C:estado_arte}

\section{El ransomware como amenaza empresarial}

\subsection{Origen y evolución histórica}

El ransomware no es una amenaza reciente. Su origen se remonta a 1989,
cuando el biólogo Joseph Popp distribuyó aproximadamente 20.000 disquetes
en la conferencia internacional sobre el SIDA de la OMS. El disquete
contenía un programa que, tras noventa reinicios del sistema, cifraba los
nombres de los ficheros y exigía el envío de 189 dólares a un apartado de
correos en Panamá. Este troyano, conocido como AIDS o PC Cyborg, utilizaba
criptografía simétrica con una clave almacenada en el propio programa, lo
que permitía recuperar los ficheros sin pagar el rescate.

La barrera que impedía la proliferación masiva del ransomware durante los
años noventa y la primera década de los dos mil era la ausencia de un
mecanismo anónimo de pago. Esta barrera desapareció con la aparición de
Bitcoin en 2009, que proporcionó por primera vez un sistema de pago
pseudoanónimo, descentralizado e irreversible.

El primer ransomware moderno fue CryptoLocker, aparecido en septiembre
de 2013. Utilizaba criptografía asimétrica RSA de 2.048~bits para proteger
la clave de cifrado AES, de modo que sin el servidor de mando y control era
matemáticamente imposible recuperar los ficheros. Infectó más de 250.000
sistemas en sus primeros cuatro meses, recaudando entre tres y veintisiete
millones de dólares.

El salto cualitativo llegó en 2017 con WannaCry, que explotaba la
vulnerabilidad EternalBlue del protocolo SMB de Windows, filtrada de los
sistemas de la \gls{nsa} por el grupo Shadow Brokers. En pocas horas
infectó más de 230.000 sistemas en 150 países, afectando al Servicio
Nacional de Salud del Reino Unido, a Telefónica y a organismos
gubernamentales de todo el mundo.

\subsection{El modelo Ransomware-as-a-Service}

La transformación más significativa del ecosistema del ransomware en los
últimos años no es técnica sino organizativa. El modelo \gls{raas} ha
convertido el ransomware en una industria estructurada con división del
trabajo. Los grupos desarrolladores mantienen la infraestructura técnica:
el cifrador, el servidor de mando y control y el portal de negociación
con las víctimas. Estos grupos ceden el acceso a afiliados que ejecutan
los ataques a cambio de una comisión de entre el 20 y el 30\,\% del
rescate cobrado.

Los grupos más activos en el período 2022-2024 son LockBit, responsable
de más del 25\,\% de los ataques conocidos; BlackCat (ALPHV), que introdujo
el uso de Rust como lenguaje del cifrador por su capacidad multiplataforma;
y Cl0p, que en 2023 explotó una vulnerabilidad en MOVEit para comprometer
más de 2.700 organizaciones simultáneamente.

\subsection{La doble y triple extorsión}

La doble extorsión, introducida por el grupo Maze en 2019, añade una
segunda amenaza al cifrado: antes de ejecutarlo, el atacante exfiltra una
copia de los datos sensibles y amenaza con publicarlos si no se paga el
rescate. Esto convierte el \textit{backup} en condición necesaria pero no
suficiente para la recuperación.

Algunos grupos han ido más allá con la triple extorsión, que añade ataques
de denegación de servicio distribuido (\gls{ddos}) o contacto directo con
clientes y reguladores para aumentar la presión.

\subsection{Evasión de EDR y ofuscación del cifrador}
\label{sec:evasion-edr}

Un factor que explica por qué la instrumentación mediante \gls{fim} y
análisis de comportamiento --- y no únicamente firmas de antivirus
tradicional --- se ha vuelto imprescindible es la sofisticación creciente
de las técnicas de evasión empleadas por los cifradores modernos. Familias
como BlackCat y LockBit~3.0 incorporan de forma habitual la desactivación
de soluciones de \textit{endpoint detection and response} mediante
controladores vulnerables firmados digitalmente (técnica conocida como
\textit{Bring Your Own Vulnerable Driver}, BYOVD), la ofuscación del binario mediante
empaquetadores personalizados que dificultan el análisis estático, y la
ejecución en memoria sin escritura a disco (\textit{fileless execution})
para evitar la detección basada en \textit{hash} de fichero.

Esta evolución es la justificación técnica de fondo para el enfoque
adoptado en este trabajo: dado que la detección basada en firmas o en la
reputación del binario es cada vez más evitable, la detección basada en
\emph{comportamiento observable a nivel de sistema de ficheros} --- el
principio en el que se apoya el módulo \gls{fim} de Wazuh --- resulta más
robusta frente a la evasión, puesto que un cifrador que modifica cientos
de ficheros en segundos produce el mismo patrón de eventos
independientemente de cómo se oculte el binario que lo ejecuta.

\subsection{Impacto económico y sectorial}

El informe de Veeam Data Protection Trends~2024~\cite{veeam2024} estima
que el coste medio de un ataque de ransomware supera los cuatro millones
de dólares, incluyendo el tiempo de inactividad, los costes de recuperación,
las sanciones regulatorias y el daño reputacional. El tiempo medio de
interrupción de la actividad es de veinte días.

El sector sanitario merece una mención especial por la gravedad de las
consecuencias. Durante la pandemia de COVID-19, el grupo Ryuk atacó
sistemáticamente hospitales en Estados Unidos, Alemania y Francia,
provocando la cancelación de cirugías urgentes y el desvío de ambulancias.

\subsection{El mercado de ciberseguros como indicador de mercado}
\label{sec:ciberseguros}

Un indicador indirecto pero revelador de la magnitud del problema es la
evolución del mercado de ciberseguros. Tras las oleadas masivas de
ataques de 2020 y 2021, las aseguradoras endurecieron notablemente los
requisitos de suscripción de pólizas que cubren incidentes de
ransomware: la autenticación multifactor, la segmentación de red y,
específicamente, la existencia de copias de seguridad \emph{inmutables o
desconectadas de la red de producción} pasaron de ser recomendaciones
opcionales a requisitos contractuales habituales para acceder a
cobertura o para evitar primas prohibitivas. Este endurecimiento del
mercado asegurador constituye una validación externa, ajena al ámbito
puramente técnico, de la tesis central de este trabajo: la industria del
riesgo ha identificado empíricamente que la disponibilidad de un backup
resistente a la manipulación por parte del atacante es el factor
diferencial que determina si un siniestro de ransomware es un incidente
recuperable o una pérdida catastrófica.

\section{Marcos de madurez de resiliencia frente a ransomware}
\label{sec:marcos-madurez}

Además de los marcos regulatorios NIS2 y ENS introducidos en el
capítulo~\ref{C:introduccion}, distintas organizaciones de estandarización han
publicado marcos de madurez específicos para la resiliencia frente a
ransomware, que resulta útil revisar para situar correctamente la contribución
de este trabajo dentro de un marco de referencia más amplio que la comparación
estrictamente técnica de dos herramientas de \emph{backup}.

El \emph{framework} de ciberseguridad del Instituto Nacional de Estándares y
Tecnología de Estados Unidos (\gls{nist}), en su versión~2.0, organiza las
capacidades de una organización en seis funciones: Gobernar, Identificar,
Proteger, Detectar, Responder y Recuperar. Este trabajo se sitúa
simultáneamente en dos de estas seis funciones: la función \emph{Proteger},
mediante el almacenamiento inmutable evaluado en las estrategias~B y~C, y la
función \emph{Detectar}, mediante la instrumentación \gls{siem} evaluada a
través de la hipótesis~H2. La función \emph{Recuperar} se evalúa
indirectamente a través del \gls{mttr}, aunque el laboratorio no cubre los
aspectos organizativos de esta función (comunicación de crisis, coordinación
con terceros) que el marco \gls{nist} contempla explícitamente.

El Centro de Seguridad de Internet (\gls{cis}), por su parte, incluye en sus
\emph{CIS Critical Security Controls} el control~11 (\emph{Data Recovery}),
que establece explícitamente la necesidad de mantener copias de seguridad
protegidas frente a modificación o eliminación no autorizada y de probar
periódicamente su restauración, en línea directa con la recomendación de
pruebas periódicas de recuperación formulada en la
sección~\ref{sec:implicaciones-practicas} del capítulo~\ref{C:discusion}.

La convergencia de estos tres marcos --- NIS2, ENS y los controles \gls{cis}
--- con la evidencia experimental de este \pfg no es casual: los tres han
sido revisados o actualizados en los últimos años precisamente en respuesta a
la evolución de las tácticas de destrucción de \emph{backups} descrita en la
sección siguiente, lo que confirma que el problema abordado experimentalmente
en este trabajo es reconocido de forma convergente tanto por la industria del
ciberseguro (sección~\ref{sec:ciberseguros}) como por los marcos normativos y
de buenas prácticas del sector.

\section{Tácticas de destrucción de backups}
\label{sec:destruccion-backups}

La destrucción de copias de seguridad como fase preparatoria del ataque
ha evolucionado desde una táctica ocasional hasta un componente sistemático
de la mayoría de los kits \gls{raas} desde 2019.

En entornos Windows, las técnicas más frecuentes son la eliminación de las
instantáneas de volumen (\textit{Shadow Copies}) mediante \texttt{vssadmin
delete shadows /all /quiet}, la desactivación del servicio de copia de
seguridad y la eliminación de los puntos de restauración del sistema.

En entornos Linux y almacenamiento en red, las técnicas equivalentes son
el borrado de repositorios Restic o Duplicati mediante credenciales robadas,
la eliminación de \textit{snapshots} en ZFS o Btrfs, y la sobrescritura de
\textit{backups} en \textit{buckets} S3 mediante claves de acceso
comprometidas.

Un aspecto frecuentemente subestimado es el período de reconocimiento
previo al cifrado. Según el informe IBM X-Force~2024~\cite{ibm2024},
el tiempo medio entre la intrusión inicial y el despliegue del cifrador
supera los 200~días. Durante ese período el atacante identifica y prepara
la destrucción de las copias de seguridad.

La magnitud del problema queda cuantificada con precisión en el informe
\emph{The State of Ransomware~2025} de Sophos~\cite{sophos2025}: el 94\,\%
de las organizaciones que sufrieron un ataque de ransomware en el último año
declaran que los atacantes intentaron comprometer explícitamente sus copias de
seguridad durante el ataque, y ese intento tuvo éxito en el 57\,\% de los
casos, es decir, en más de la mitad de las organizaciones afectadas. El mismo
informe documenta una caída hasta un mínimo de cuatro años en el uso efectivo de
\emph{backups} como vía de recuperación (53\,\%, frente al 73\,\% del año
anterior), un dato que sugiere que la erosión de la confianza en las estrategias
de \emph{backup} convencionales --- precisamente la motivación central de este
\pfg --- no es una percepción anecdótica sino una tendencia sectorial medible.
Estos datos, publicados con posterioridad al diseño del experimento pero antes
de la redacción de esta memoria, refuerzan \emph{a posteriori} la relevancia
del escenario~E2 como el más representativo del comportamiento real de los
atacantes, y no como un caso extremo poco frecuente.

\section{Estrategias de backup y almacenamiento inmutable}

\subsection{La regla 3-2-1 y sus extensiones}

La regla 3-2-1 de \textit{backup}, popularizada por Peter Krogh en
2005~\cite{krogh2005}, establece que deben mantenerse tres copias de los
datos, en dos medios distintos, con al menos una copia fuera de las
instalaciones. Esta regla proporciona protección contra fallos de hardware
y desastres físicos, pero no contempla el escenario de un atacante con
acceso administrativo que puede operar sobre todas las copias
simultáneamente.

La extensión más adoptada es la regla 3-2-1-1-0, que añade: al menos una
copia debe estar \textit{offline} o en un medio inmutable, y debe
verificarse que hay cero errores mediante pruebas de restauración periódicas.

\subsection{Restic como herramienta de backup}

Restic~\cite{restic2024} es una herramienta de \textit{backup} de código
abierto escrita en Go que cifra todos los datos antes de almacenarlos
usando AES-256-CTR con autenticación Poly1305. Su diseño basado en
deduplicación a nivel de bloque lo hace eficiente para \textit{backups}
incrementales. Soporta múltiples \textit{backends} incluyendo \gls{sftp},
Amazon \gls{s3} y servicios compatibles como MinIO.

Una consideración importante en el uso de Restic con Object Lock es la
gestión de los ficheros de bloqueo del repositorio. Con Object Lock
COMPLIANCE activo, Restic no puede eliminar estos \textit{locks} al
completar una operación. La solución es el uso del flag
\texttt{--no-lock} en las operaciones de \textit{backup} y restauración.

\subsection{S3 Object Lock y almacenamiento inmutable}

S3 Object Lock~\cite{aws2023} es una característica introducida por Amazon
en 2018 que permite proteger los objetos almacenados frente a modificaciones
y eliminaciones durante un período definido. Implementa el modelo WORM
(\gls{worm}).

Object Lock admite dos modos de retención:

\begin{description}
  \item[Governance] Protege los objetos frente a la mayoría de los usuarios,
  pero permite que usuarios con permisos especiales anulen la protección.

  \item[Compliance] Proporciona una protección absoluta e irrevocable: ningún
  usuario, incluyendo el usuario raíz, puede eliminar un objeto protegido
  antes de que expire su período de retención. Este modo se utiliza en la
  estrategia~C de este trabajo.
\end{description}

MinIO~\cite{minio2024} es una implementación de código abierto de un
servidor de almacenamiento de objetos compatible con la \gls{api} de
Amazon \gls{s3}. Implementa la especificación completa de S3 incluyendo
Object Lock, versionado y políticas \gls{iam}, lo que lo hace adecuado
para despliegues \textit{on-premise} sin dependencia de servicios cloud.

\subsection{Alternativas al modelo Object Lock}

S3 Object Lock no es la única vía para conseguir almacenamiento
\gls{worm}, y conviene situarla frente a las alternativas para justificar
su elección en este trabajo.

\begin{description}
  \item[Cinta LTO con \textit{write-protect}] La cinta magnética LTO
  dispone de una pestaña física de protección de escritura que la hace
  inmune por diseño a cualquier ataque puramente lógico: un ransomware no
  puede sobrescribir una cinta cuya pestaña esté en posición de solo
  lectura, con independencia de los privilegios obtenidos en la red. Su
  inconveniente principal es operativo: requiere intervención humana para
  rotar y proteger físicamente cada cinta, lo que introduce demoras del
  orden de horas o días en el \gls{rto} y no encaja con estrategias de
  recuperación que exigen automatización completa.

  \item[\textit{Snapshots} de ZFS y Btrfs con \texttt{send}/\texttt{receive}]
  Ambos sistemas de ficheros permiten crear instantáneas de solo lectura y
  replicarlas hacia un segundo host mediante un flujo unidireccional
  (\texttt{zfs send \textbar{} zfs receive}). Si el host de destino no
  expone ningún servicio de administración remota alcanzable desde el
  origen, un atacante que compromete el fileserver no puede alcanzar ni
  destruir las instantáneas remotas. La limitación frente a Object Lock es
  que la protección depende de una correcta segmentación de red: un fallo
  de configuración que exponga el host de réplicas anula la garantía.

  \item[\textit{Immutable Blob Storage} de Azure y \textit{Bucket Lock} de
  Google Cloud] Ambos proveedores ofrecen mecanismos funcionalmente
  equivalentes a S3 Object Lock COMPLIANCE: una vez fijada la política de
  retención, ni siquiera el propietario de la suscripción puede eliminar
  el objeto antes de que expire el plazo. La elección de MinIO en este
  trabajo, frente a estas alternativas cloud, responde a la necesidad de
  un entorno completamente \textit{on-premise} y reproducible sin
  dependencia de credenciales de terceros, requisito imprescindible para
  la reproducibilidad del laboratorio descrita en el anexo~C.
\end{description}

Las tres alternativas comparten con Object Lock el principio de fondo que
este trabajo pone a prueba experimentalmente: la inmutabilidad debe ser
\emph{estructural} (impuesta por el sistema, no revocable por ningún
credential) y no meramente \emph{procedimental} (basada en permisos que un
atacante con privilegios suficientes puede modificar).

\section{Sistemas SIEM y detección de ransomware}

\subsection{Arquitectura de un SIEM}

Un \gls{siem} combina la gestión de información de seguridad y la gestión
de eventos de seguridad en una plataforma unificada. Su arquitectura
típica se organiza en cuatro capas: recopilación de eventos de múltiples
fuentes, normalización a un formato común, correlación mediante reglas y
modelos de comportamiento, y presentación mediante cuadros de mando y
herramientas de investigación forense.

\subsection{Wazuh como plataforma SIEM y XDR}

Wazuh~\cite{wazuh2024} es una plataforma \gls{siem} y \gls{xdr} de código
abierto derivada de OSSEC. Su arquitectura se compone de tres elementos:

\begin{itemize}
  \item \textbf{Agente}: proceso ligero desplegado en cada sistema
  monitorizado. Recoge eventos del sistema operativo, monitoriza la
  integridad de ficheros y envía todo al manager mediante \gls{tls} mutuo.
  Consume menos de 35~MB de memoria en reposo.

  \item \textbf{Manager}: núcleo del sistema. Decodifica los eventos
  mediante expresiones regulares y XML, aplica las reglas de correlación
  y genera las alertas.

  \item \textbf{Indexer}: basado en OpenSearch, almacena todas las alertas
  y permite su consulta mediante una \gls{api} REST.
\end{itemize}

\subsection{File Integrity Monitoring en Wazuh}

El \gls{fim} es la capacidad más relevante de Wazuh para la detección de
ransomware. El módulo \texttt{syscheck} monitoriza directorios y ficheros
para detectar cualquier cambio: creación, modificación, eliminación o
cambio de permisos. En modo \textit{realtime}, utiliza \texttt{inotify}
en Linux para detectar cambios inmediatamente sin esperar a un ciclo de
escaneo periódico.

La detección de ransomware mediante \gls{fim} se basa en la observación de
que un cifrador masivo produce un número anómalo de eventos de modificación
o creación de ficheros en un período breve. Las reglas de frecuencia de
Wazuh permiten configurar umbrales: si se detectan más de N cambios en T
segundos, se genera una alerta de nivel elevado.

\subsection{Alternativas SIEM de código abierto}

La elección de Wazuh frente a otras plataformas SIEM de código abierto
merece justificación explícita, dado que el ecosistema ofrece varias
alternativas maduras.

\begin{description}
  \item[OSSEC clásico] Wazuh es históricamente un \textit{fork} de OSSEC,
  del que hereda el motor de reglas y el módulo de \gls{fim}. OSSEC
  clásico permanece disponible pero su desarrollo es mucho menos activo y
  carece de un componente de indexación y visualización integrado
  equivalente al Indexer/OpenSearch de Wazuh, lo que obligaría a integrar
  manualmente una solución de almacenamiento y consulta separada.

  \item[Security Onion] Distribución Linux completa orientada a
  \gls{siem} y detección de intrusiones de red, que integra Suricata, Zeek
  y una pila Elastic completa. Ofrece una cobertura de red muy superior a
  Wazuh de fábrica, pero su despliegue está pensado para una máquina o
  clúster dedicado con requisitos de recursos considerablemente mayores,
  lo que la hace menos adecuada para un laboratorio basado en contenedores
  ligeros ejecutado sobre un único equipo de sobremesa.

  \item[Suricata + pila ELK autogestionada] Permite un control total sobre
  cada componente, pero exige construir manualmente la correlación de
  reglas, el \gls{fim} y el enriquecimiento de eventos que Wazuh
  proporciona de forma nativa. El coste de ingeniería para alcanzar la
  funcionalidad de detección requerida en este trabajo sería
  considerablemente mayor que el de configurar reglas \textit{custom}
  sobre una plataforma ya integrada.
\end{description}

Wazuh se seleccionó por tres motivos concretos: dispone de un módulo
\gls{fim} en tiempo real maduro y bien documentado, que es la capacidad
central que requiere este trabajo; su arquitectura de agente ligero
(inferior a 35\,MB de memoria) encaja con el presupuesto de recursos de un
laboratorio en contenedores; y su desarrollo activo como plataforma
\gls{xdr} lo mantiene alineado con las necesidades de detección de
amenazas modernas más allá del ámbito puramente SIEM.

\subsection{Detección basada en auditd}

Una alternativa al \gls{fim} es \texttt{auditd}, que intercepta las
llamadas al sistema en el momento en que se producen. Proporciona mayor
granularidad pero genera un volumen muy elevado de eventos y requiere
soporte del kernel. En el laboratorio de este trabajo, \texttt{auditd} no
pudo utilizarse debido a una incompatibilidad con el kernel CachyOS del
sistema \textit{host}, documentada en el capítulo~\ref{C:arquitectura}.

\section{El framework MITRE ATT\&CK}

MITRE ATT\&CK~\cite{mitre2023} es una base de conocimiento desarrollada
por la corporación MITRE que cataloga las tácticas y técnicas empleadas
por actores de amenaza reales. La matriz ATT\&CK para Enterprise organiza
más de 600~técnicas y sub-técnicas en catorce tácticas que abarcan el
ciclo de vida completo de un ataque.

Los cuatro escenarios de este trabajo mapean directamente con técnicas de
la matriz:

\begin{table}[!h]
  \centering
  \caption[Correspondencia entre escenarios del laboratorio y técnicas
           MITRE ATT\&CK.]
          {Correspondencia entre escenarios del laboratorio y técnicas
           MITRE ATT\&CK.}
  \label{tab:attck}
  \begin{tabularx}{\textwidth}{C{1.8cm} C{2.8cm} X C{2.5cm}}
    \toprule
    \rowcolor{tabla-cabecera}
    \textcolor{white}{\textbf{Escenario}} &
    \textcolor{white}{\textbf{ID ATT\&CK}} &
    \textcolor{white}{\textbf{Nombre}} &
    \textcolor{white}{\textbf{Táctica}} \\
    \midrule
    E1 & T1486      & Data Encrypted for Impact  & Impacto \\
    E2 & T1490      & Inhibit System Recovery     & Impacto \\
    E3 & T1021.004  & Remote Services: SSH        & Mov. lateral \\
    E4 & T1041      & Exfiltration Over C2        & Exfiltración \\
    \bottomrule
  \end{tabularx}
\end{table}

Atomic Red Team~\cite{atomicredteam2024} es un \textit{framework} de
simulación de ataques que proporciona implementaciones de código abierto
de las técnicas ATT\&CK. Sus tests están diseñados para ser atómicos, es
decir, ejecutan una única técnica en aislamiento y son ampliamente
utilizados para validar la cobertura de herramientas de detección.

\section{Trabajos relacionados}

La literatura académica sobre evaluación experimental de estrategias de
backup en escenarios de ransomware es escasa.

Moussaileb et al.~\cite{moussaileb2021} proponen un sistema de detección basado en el
análisis del acceso a ficheros en tiempo real, con tasas de verdaderos
positivos superiores al 95\,\% en condiciones de laboratorio.

Al-Mashhadani y Shafiq~\cite{almashhadani2019} evalúan técnicas de detección mediante
aprendizaje automático sobre tráfico de red, logrando tasas de detección
superiores al 98\,\% con tiempos inferiores a 30~segundos para el
ransomware Locky.

Cabaj y Mazurczyk~\cite{cabaj2016} analizan el comportamiento de red del ransomware
CryptoWall y proponen un sistema de detección y mitigación basado en
\gls{sdn}.

Fuera de los tres trabajos centrados en detección, tres líneas de investigación
recientes resultan particularmente relevantes para situar metodológicamente este
\pfg, aunque abordan preguntas distintas a la comparación de estrategias de
recuperación.

Gehani et al.~\cite{gehani2023} abordan un problema distinto pero
complementario: la falta de conjuntos de datos abiertos y reproducibles sobre el
comportamiento real de binarios de ransomware. Su trabajo construye un
\emph{dataset} de \emph{logs} de actividad y grafos de procedencia a partir de la
ejecución de muestras reales en \emph{sandbox}, extraídas de repositorios como
MalwareBazaar, y publica tanto los datos como el código de orquestación necesario
para regenerarlos. Comparte con este trabajo el énfasis en la reproducibilidad
como requisito metodológico de primer orden --- de hecho, el título del artículo,
\emph{Towards Reproducible Ransomware Analysis}, podría aplicarse igualmente al
laboratorio aquí presentado --- pero se centra en la caracterización del
comportamiento del cifrador en sí, no en la comparación de estrategias de
recuperación posteriores al ataque.

Von der Assen et al.~\cite{vonderassen2024} proponen GuardFS, un sistema de
ficheros \emph{overlay} para Linux que integra detección y mitigación de
ransomware en la misma capa: intercepta el acceso a ficheros, extrae
características del patrón de acceso mediante modelos entrenados y aplica una de
tres estrategias de defensa activa (ofuscación, retardo o seguimiento del acceso)
en tiempo real, antes de que se complete el cifrado. Este enfoque es
complementario al de este \pfg: mientras que GuardFS actúa en la capa del sistema
de ficheros para \emph{prevenir} o retrasar la pérdida de datos durante el
ataque, el presente trabajo asume que el cifrado no puede evitarse por completo y
se centra en garantizar la \emph{recuperación} posterior mediante almacenamiento
inmutable, junto con una detección basada en \gls{fim} a nivel de \gls{siem} en
lugar de a nivel de sistema de ficheros. Los resultados de GuardFS, que muestran
que la pérdida de datos se reduce pero no se elimina completamente incluso con
mitigación activa a nivel de \emph{filesystem}, refuerzan indirectamente la tesis
de este trabajo: ninguna capa de defensa preventiva es suficiente por sí sola, y
la garantía última de recuperación --- el objeto de estudio de este \pfg ---
sigue siendo necesaria como última línea de defensa.

Bitzki et al.~\cite{netsecbed2026} presentan NetSecBed, un \emph{testbed}
nativo en contenedores para la generación reproducible de evidencias de
experimentación en ciberseguridad, con 60~escenarios de ataque y nueve servicios
objetivo implementados como contenedores de propósito único. Aunque orientado a
entornos IoT/IIoT y no específicamente al ransomware, NetSecBed valida
metodológicamente la decisión arquitectónica adoptada en el capítulo~\ref{C:arquitectura}
de este trabajo: el uso de contenedores de propósito único orquestados
declarativamente como mecanismo para lograr reproducibilidad, extensibilidad y
trazabilidad en experimentación de ciberseguridad es una tendencia consolidada en
la literatura reciente, y no una decisión de conveniencia aislada de este \pfg.

Ninguno de estos trabajos aborda de forma experimental y controlada la
comparación entre estrategias de backup en términos de \gls{mttr}, ni
incluye la destrucción activa de las copias como variable experimental.
El presente trabajo cubre ese vacío mediante un diseño factorial
completamente automatizado y con análisis estadístico riguroso.

\subsection{Análisis comparativo del vacío experimental}

Se incluyen adicionalmente en la comparación los tres trabajos metodológicamente
próximos discutidos anteriormente. Aunque Gehani et al.~\cite{gehani2023}
alcanzan un grado de automatización y reproducibilidad comparable al de este
\pfg (columna «Automatizado y reproducible»), su objeto de estudio es la propia
ejecución del ransomware, no la comparación de estrategias de recuperación
posteriores. Von der Assen et al.~\cite{vonderassen2024} evalúan
sistemáticamente distintas configuraciones de defensa con métricas cuantitativas,
lo que se refleja en la columna «Diseño repetido + estadística» como cobertura
parcial, pero no comparan estrategias de \emph{backup} ni modelan la destrucción
activa de copias. Bitzki et al.~\cite{netsecbed2026} logran un alto grado de
automatización mediante contenedores de propósito único, en línea con la
arquitectura de este trabajo, pero su dominio de aplicación (tráfico de red
IoT/IIoT) es ajeno al ransomware y a la evaluación de estrategias de recuperación.

La Tabla~\ref{tab:gap-analysis} sintetiza las dimensiones metodológicas
cubiertas por cada trabajo relacionado frente a las cubiertas por el
presente \pfg. Se han seleccionado cuatro dimensiones que resultan
particularmente relevantes para el objetivo de este trabajo: si el estudio
incluye la destrucción activa de copias de seguridad como parte del
escenario de amenaza, si compara estrategias de recuperación (y no
únicamente mecanismos de detección), si aplica un diseño experimental
repetido con análisis estadístico de significancia, y si el entorno es
completamente automatizado y reproducible por terceros.

\begin{table}[!h]
  \centering
  \caption[Comparación metodológica con los trabajos relacionados.]
          {Comparación metodológica con los trabajos relacionados.}
  \label{tab:gap-analysis}
  \begin{tabularx}{\textwidth}{X C{2.6cm} C{2.6cm} C{2.6cm} C{2.6cm}}
    \toprule
    \rowcolor{tabla-cabecera}
    \textcolor{white}{\textbf{Trabajo}} &
    \textcolor{white}{\textbf{Destrucción de backups}} &
    \textcolor{white}{\textbf{Compara recuperación}} &
    \textcolor{white}{\textbf{Diseño repetido + estadística}} &
    \textcolor{white}{\textbf{Automatizado y reproducible}} \\
    \midrule
    Moussaileb et al.~\cite{moussaileb2021}      & No & No & No & No \\
    Al-Mashhadani y Shafiq~\cite{almashhadani2019} & No & No & Parcial & No \\
    Cabaj y Mazurczyk~\cite{cabaj2016}            & No & No & No & No \\
    Gehani et al.~\cite{gehani2023}                & No & No & No & Parcial \\
    Von der Assen et al.~\cite{vonderassen2024}    & No & No & Parcial & No \\
    Bitzki et al.~\cite{netsecbed2026}              & No & No & No & Sí \\
    \textbf{Este trabajo}                         & \textbf{Sí} & \textbf{Sí} & \textbf{Sí} & \textbf{Sí} \\
    \bottomrule
  \end{tabularx}
\end{table}

Los tres trabajos revisados comparten un enfoque centrado exclusivamente
en la \emph{detección} del ransomware --- mediante acceso a ficheros,
tráfico de red o ambos --- evaluada sobre un único conjunto de muestras
sin repeticiones ni análisis de significancia estadística. Ninguno modela
la fase de destrucción de copias de seguridad como parte del escenario de
amenaza, pese a que, como se ha argumentado en la
sección~\ref{sec:destruccion-backups}, esta fase es un
componente sistemático del ransomware moderno desde 2019. Esta ausencia es
precisamente el vacío que motiva el diseño experimental del
capítulo~\ref{C:disenyo}: sin una variable experimental que capture la
destrucción activa del repositorio, ningún estudio previo puede
diferenciar entre una estrategia de backup rápida y una estrategia de
backup \emph{disponible}, distinción que resulta central en los
resultados de este trabajo (capítulo~\ref{C:resultados}).

\end{document}
